\section{DISCUSSÃO E LIMITAÇÕES}

\subsection{Os resultados à luz da literatura}

Os dados reproduzem o padrão de compressão descrito para o Airbnb \citep{zervas2021}: 32,4\% dos anúncios elegíveis têm nota exatamente igual a 5,0. A maior parte das sentenças é positiva em ambos os modelos de sentimento (78,9\% em A e 81,5\% em B). O resultado concorda com a observação de que as avaliações da plataforma adotam uma norma de comentário positivo \citep{bridges2018} e com a comparação, em cidades brasileiras e americanas, que encontra avaliações do Airbnb mais positivas que as de uma plataforma hoteleira \citep{santos2020}. Dado esse viés, a decisão de trabalhar com a ordenação do sentimento e não com o seu nível e de tomar o percentil da nota como alvo, tem apoio nos próprios dados.

A decomposição da nota coincide, em parte, com os temas que a literatura aponta para a experiência do hóspede. \citet{cheng2019} destacam localização, comodidades e anfitrião. \citet{sutherland2020} agrupam os tópicos em localização, características do imóvel e serviço do anfitrião. Neste estudo, comodidades e limpeza respondem pela maior parte da variação explicada da nota global. As queixas de comunicação com o anfitrião também se associam a notas menores. A localização, por outro lado, quase não contribui para a nota global, embora seja o aspecto que mais explica a sua própria subnota. A diferença em relação aos estudos citados está no desfecho: eles descrevem os temas mais frequentes ou mais valorados nas avaliações. Este trabalho mede a associação de cada aspecto com a nota numérica. Os dois resultados são compatíveis. A localização aparece em muitas avaliações e o hóspede a avalia em separado, de modo que é plausível que a sua variação entre anúncios acrescente pouco à nota global depois de considerados os atributos do imóvel, hipótese que os dados não testam. A análise de \citet{luo2019}, que relaciona aspectos latentes à nota numérica, é o antecedente mais próximo. Este estudo a estende com validação da medida por aspecto, dois modelos de sentimento e um corpus multilíngue.

A leitura por tópico traz uma assimetria que a teoria dos três fatores (Seção 2.1) ajuda a descrever. As queixas da Tabela~
ef{tab:queixas} se associam a diferenças de 10 a 26 pontos no percentil da nota e os elogios da Tabela~
ef{tab:elogios}, a diferenças de 5 a 11 pontos. O contraste é mais nítido em limpeza: a queixa de sujeira tem a maior diferença associada da Tabela~
ef{tab:queixas} ($-26{,}4$ pontos) e o elogio de limpeza, a menor da Tabela~
ef{tab:elogios} ($+4{,}8$). As queixas de instalações e manutenção (banheiro, ar-condicionado, cozinha, cama, conexão de internet e reparos) e de cheiro estão entre $-17{,}7$ e $-21{,}8$ pontos. Esse padrão é compatível com atributos que se comportam como fatores básicos: o mau desempenho penaliza a nota mais do que o bom desempenho a recompensa. Limpeza e as instalações do imóvel seriam os candidatos. Os elogios de maior diferença associada são os de satisfação geral (sentir-se em casa, $+10{,}9$ pontos) e de comunicação com o anfitrião ($+8{,}5$), mais próximos de fatores de desempenho ou de entusiasmo. A leitura é descritiva e tem limites. Os tópicos comparados são diferentes e muito desiguais em prevalência (o elogio ao anfitrião aparece em 84\% dos anúncios e a queixa de sujeira em 7\%), as menções vêm de previsões com erro e os aspectos de comodidades e de satisfação geral estão na classe de cautela. A assimetria entre queixa e elogio de um mesmo atributo não foi testada formalmente. Esse teste é um dos desdobramentos indicados nas Conclusões.

O resultado nulo do índice sintético tem relação com o efeito halo \citep{thorndike1920}. Como as subnotas e a nota global são muito correlacionadas entre si, a nota oficial já captura boa parte do que o sentimento médio dos aspectos poderia acrescentar. O índice reordena os anúncios de nota 5,0 em relação à nota oficial, mas essa reordenação não se associa ao volume de avaliações além do que a nota oficial já explica. Os desfechos disponíveis, que são índices de volume, oferecem um teste pouco favorável ao índice. Essa é a primeira das duas leituras do resultado apresentadas na Seção 7.

\subsection{Implicações para anfitriões e empresas gestoras}

O resultado mais utilizável está nas queixas por tópico. Entre as 15 queixas de maior peso, 13 se repetem nos dois modelos de sentimento. Entre elas figuram banheiro, ar-condicionado, cozinha, cama, conexão de internet, manutenção, limpeza, cheiro e resposta do anfitrião. As queixas se associam a quedas de 10 a 26 pontos no percentil da nota. O diagnóstico por anúncio permite localizar, entre essas queixas, as que se repetem em um anúncio específico. O instrumento é descritivo. A associação é uma média da cidade. Uma queixa de banheiro em um anúncio não implica que a correção do banheiro eleve a sua nota.

\subsection{Limitações}

O sentimento foi calculado sobre as 30 avaliações mais recentes de cada anúncio. O limite atinge 32\% dos anúncios elegíveis (9.266), embora apenas 5,2\% dos de nota 5,0 e reduz a confiabilidade do sentimento médio por anúncio, cuja correção de Spearman-Brown é de 0,50 com 30 avaliações. Os dados contêm apenas anúncios ativos na coleta. A data da avaliação não indica a data da estadia.

Não há ocupação nem receita observadas. Os campos disponíveis são funções do número de avaliações recentes. Nenhum resultado diz respeito a desempenho econômico.

Os rótulos de referência vêm de um único pesquisador, que também revisou o mapa dos tópicos para aspectos. Um segundo anotador rotulou 100 sentenças, com concordância alta no sentimento (kappa de 0,86) e moderada no aspecto (kappa de 0,62), de modo que a atribuição de aspecto carrega erro de rotulagem além do erro do modelo. A amostra rotulada tem 297 sentenças e, ponderada pela população, tamanho efetivo de cerca de 72, razão pela qual os intervalos são largos e a comparação entre modelos é frágil: a diferença de kappa entre B e A é de $-0{,}141$, com intervalo que mal exclui zero. A confiabilidade das metades do sentimento por aspecto é baixa; as subnotas usadas como referência são comprimidas. Sete dos dez aspectos pedem cautela. A classificação dos aspectos é indicativa: com 27 sentenças rotuladas por aspecto, os intervalos de Wilson são largos e os dos três aspectos confiáveis incluem o limiar de 0,60 e em limpeza os dois anotadores concordam pouco (33\%, $n = 6$). Comodidades, o aspecto de maior peso, tem atribuição imprecisa (44\%, intervalo de Wilson de 28\% a 63\%) e não tem subnota que o valide; dividido em subaspectos, o seu peso se reparte sem que um deles domine (Seção 5). O seu conteúdo se sobrepõe em parte ao de outros aspectos. O LDA permite checar apenas quatro aspectos.

Os modelos de sentimento divergem de forma que a escolha não é óbvia. O modelo A é o principal por uma regra fixada antes da execução do B, mas B é melhor na amostra rotulada e A, na população. Por isso os resultados de B acompanham todos os blocos. Onde os dois divergem (por exemplo, nos coeficientes dentro da amostra do índice sintético), as duas leituras são relatadas sem preferência.

A importância relativa de um aspecto mistura efeito e confiabilidade da medida; as associações com a nota são entre anúncios e não dentro de um mesmo anúncio. O texto e a nota vêm dos mesmos hóspedes, de modo que a análise decompõe o conteúdo da nota e não a prevê por um fator independente e o texto cobre as 30 avaliações mais recentes enquanto a nota é acumulada. Anúncios sem aspectos suficientes têm o sentimento imputado. Quanto ao índice sintético, apenas 5.142 dos anúncios de nota 5,0 o possuem. Os que ficam de fora têm menos avaliações, de modo que as conclusões sobre o teto valem para esse subconjunto. A decisão de não aplicar encolhimento foi tomada depois de ver o resultado com encolhimento, o que justifica reportar as duas versões. Os coeficientes dentro da amostra não são robustos e não foram confirmados fora dela.

A análise cobre um município, uma plataforma e uma coleta e os resultados podem não se estender a outros mercados.
